\documentclass[10pt,a4paper]{ctexart}

\usepackage[margin=1.05cm]{geometry}
\usepackage{enumitem}
\usepackage{titlesec}
\usepackage{hyperref}
\usepackage{xcolor}
\usepackage{setspace}

\pagestyle{empty}
\hypersetup{hidelinks}
\setlength{\parindent}{0pt}
\setlength{\parskip}{2.5pt}
\setstretch{1.02}

\setlist[itemize]{leftmargin=1.2em,itemsep=0.5pt,topsep=1pt}
\setlist[enumerate]{leftmargin=1.4em,itemsep=0.5pt,topsep=1pt}

\titleformat{\section}{\large\bfseries}{\thesection}{0.5em}{}
\titleformat{\subsection}{\normalsize\bfseries}{\thesubsection}{0.5em}{}
\titlespacing*{\section}{0pt}{5pt}{2pt}
\titlespacing*{\subsection}{0pt}{3pt}{1pt}

\begin{document}
\thispagestyle{empty}

\begin{center}
{\LARGE \textbf{金融求职定制化 AI 模拟面试助手Mock-Me Product Memo}}
\end{center}

\section{目标用户与核心痛点}

本项目面向准备金融岗位面试的学生，重点覆盖一级市场、二级市场和量化/金融科技方向。

调研投行、行研、量化方向同学后，我发现商科面试最难训练的不是自我介绍、行为面八股或纯知识题，而是\textbf{围绕个人简历、目标岗位和业务场景展开的动态追问}。典型痛点包括：简历经历经不起深挖；课程、工具、项目方法写了但答不细；项目表达详略失衡；面对公司业务 case 只能泛泛而谈；压力面下表达节奏崩；面试后缺少具体反馈。

因此，一个有价值的 AI 面试产品不应只是随机抽题，而应围绕用户材料进行定制化提问，并把面试结果沉淀为可复盘的个人训练资产。

\section{产品设计说明}

\subsection{产品定位}

本项目定位为面向金融/商科求职的\textbf{定制化 AI 模拟面试助手}。它不是通用聊天机器人，也不是题库工具，而是围绕“候选人本人 + 目标岗位 + 真实材料”展开的一对一训练系统。

核心目标是让 AI 更像一个了解候选人背景和岗位要求的面试官：基于简历、JD、作品提问；结合不同金融岗位调整关注点；根据回答继续追问或降阶引导；面试后输出结构化反馈和知识卡片。

\subsection{为什么选择这个方向}

商科求职准备中，通用资料很多，但高质量的一对一反馈很少。现有工具多偏向通用问答、英文口语、行为面题库或面经检索，能回答“可能问什么”，但很难回答“我这段经历会被怎么追问”“我投这个岗位哪里最容易露怯”。

真实金融面试通常验证三件事：是否真的做过、是否真的理解、是否真的匹配。投行、行研、量化的追问路径不同；即使岗位相同，不同简历也会导向不同问题。因此我没有把核心做成“知识库 + RAG + 题库抽题”，而是聚焦材料和岗位驱动的定制化面试。

产品闭环是：\textbf{上传简历/作品/JD → 生成定制化面试 → 实时追问与澄清 → 结构化评分 → 知识卡片与历史记录}。

\subsection{当前版本核心功能}

当前版本围绕一场定制化面试的完整流程展开。

\textbf{面试前：对齐人和岗位。} 用户选择金融/商科方向，并上传简历、作品或粘贴 JD。系统据此获取候选人经历、课程、技能、目标岗位、JD 要求和作品材料。

当前支持四种模式：简历深挖、业务模拟、混合面试、作品答辩。简历深挖只围绕简历连续追问；业务模拟基于 JD 生成岗位场景；混合面试综合简历、业务、知识、行为和压力追问；作品答辩围绕报告、PPT、投资备忘录等材料追问假设、证据链和结论稳健性。

\textbf{面试中：逼近真实互动。} 当前版本支持压力面指数、向面试官澄清和复合问题生成。压力面指数控制追问密度和质疑概率；澄清机制允许候选人先确认题意且不单独计分；复合问题把简历、业务、技术、行为面打通，更接近真实面试官提问。

\textbf{面试后：形成复盘资产。} 系统生成结构化报告，包括总分、逐题评分、优势、改进方向、遗漏知识点和知识卡片。知识卡片保存到本机历史库；报告页只展示模型精选出的 6 个重点卡片，避免信息过载。

支撑体验包括中文材料解析、整场限时、每题计时、语音回答和本机历史记录。它们都服务于同一目标：让用户围绕自己的材料和岗位反复训练，而不是刷通用题。

\subsection{取舍}

在 16 小时限制下，我主动收窄边界：不做复杂视频面试，不做海量固定题库，不做泛行业平台，不堆 UI 特效。优先保证“输入材料 → 定制出题 → 追问互动 → 反馈复盘”的最小闭环可用。

\section{版本迭代记录}

\subsection{最初方案}

项目最初基于开源 mock interviewer项目https://github.com/koushik-gupta/Composable-AI-Mock-Interviewer- 改造，原始项目基于计算机场景且定制化作用非常弱，我对原始项目进行了大刀阔斧的改进，先搭起可运行 demo，再替换行业设定、模式分类、问题生成逻辑和反馈体系。

\subsection{核心问题}

问题主要分为内容层和使用形式层。内容层更关键，因为产品差异化来自“问得准不准、像不像真实金融面试”。

\textbf{内容层问题}包括：问题太泛，像通用题库；业务题太大，像 research memo；问题压缩后上下文不足；模型可能混淆系统提示和用户回答；候选人不会时机械重复追问；不同模式边界不清；知识卡片太多、同质化。

\textbf{使用形式问题}包括：文字输入不够接近真实口头面试；面试页曾有冗余模块；难度、压力面、计时等前端表达不统一；刷新、历史保存、报告完整性影响连续训练；中文显示和 Markdown/公式渲染需要修复。

\subsection{关键修改}

\textbf{面向内容质量：} 约束问题必须可现场回答；限制每题只考察一个核心判断；要求业务题提供角色、目标、约束和假设；增加经历追问模板；简历深挖中加入专业课和工具追问；明确四种模式边界；加入复合问题、压力面指数和降阶引导；对知识卡片去重并精选展示。

\textbf{面向使用形式：} 增加澄清机制、限时和计时；加入语音回答；去掉冗余输入区和答题提示；统一前端文案；增加本机历史记录和知识库入口。

这些迭代把系统从“会出题的聊天工具”推进到“可训练、可复盘、可持续迭代的定制化金融面试系统”。

\section{产品价值}

相比直接使用通用大模型，本产品的增量在于：

\begin{itemize}
  \item 基于简历、JD、作品进行个性化提问；
  \item 连续追问，逼近真实面试的压力和层次；
  \item 融合简历、业务、知识、行为面，而不是割裂成题库；
  \item 面试后给出逐题反馈和遗漏知识点；
  \item 把知识卡片和历史记录沉淀为个人训练资产。
\end{itemize}

关键不在于多问几道题，而在于提供一个能针对自己材料持续追问并指出盲点的训练环境。

\section{下一步设计}

下一步会沿内容、形式、技术三条线推进。

\textbf{内容层：} 完善评分体系，细化回答长度、应答时间、知识正确性、结构化表达和岗位匹配度；优化知识卡片，支持按行业、岗位、课程、经历筛选；强化问题生成，区分不同模式和金融岗位；增强历史复盘；细化投行 VP、买方研究员、量化带教等面试官 persona。

\textbf{形式层：} 在语音回答基础上增加口头表达复盘，分析语速、停顿、重复表达和结构；补充开场、自我介绍、追问、结束反馈等阶段感，让体验更像完整 mock interview。

\textbf{技术层：} 将报告生成、卡片抽取、语音转写等慢任务异步化，避免高并发阻塞；增加限流、缓存和降级机制，保证公网演示和多人访问时稳定。

\section{AI 工具使用}

本项目使用 AI 工具完成调研整理、产品定义、提示词工程、代码改造、部署调试和 bug 修复。

\begin{itemize}
  \item \textbf{Codex / AI Coding 工具：}用于代码阅读、功能实现、部署联调、bug 修复和 GitHub 提交。
  \item \textbf{DeepSeek API：}作为核心对话模型；当前使用 deepseek-v4-flash，是出于能力、成本和响应速度的平衡考虑，适合高频、多轮模拟面试。
  \item \textbf{火山语音识别 API：}支持语音回答转写。
  \item \textbf{Doubao API：}用于生成测试回答和辅助验证。
\end{itemize}

我的使用方式是：先明确目标用户和产品闭环，再让 AI 加速实现，而不是先堆能力再寻找场景。

\section{总结}

这个项目不试图成为“大而全”的 AI 面试平台，而是聚焦金融/商科求职中最难练、最依赖真人经验的一环：围绕个人材料和目标岗位展开的定制化追问。

在有限时间里，我完成了一个可公网访问、可上传真实材料、可进行多模式面试、可生成结构化反馈并沉淀知识卡片的最小闭环。它的核心价值不是“能聊天”，而是更像一个懂金融面试、能针对候选人本人持续训练的模拟面试官。

\end{document}
